\documentclass[12pt,a4paper]{article}

\usepackage[utf8]{inputenc}
\usepackage[T1]{fontenc}
\usepackage[spanish]{babel}
\usepackage{geometry}
\usepackage{xcolor}
\usepackage{listings}
\usepackage{hyperref}

\geometry{margin=2.5cm}

\definecolor{codebg}{RGB}{248,248,248}
\definecolor{codeframe}{RGB}{220,220,220}
\definecolor{codekeyword}{RGB}{0,92,175}
\definecolor{codestring}{RGB}{120,70,0}

\lstdefinestyle{pythonstyle}{
    language=Python,
    backgroundcolor=\color{codebg},
    frame=single,
    rulecolor=\color{codeframe},
    basicstyle=\ttfamily\small,
    keywordstyle=\color{codekeyword}\bfseries,
    stringstyle=\color{codestring},
    showstringspaces=false,
    breaklines=true,
    tabsize=4
}

\lstdefinestyle{textstyle}{
    backgroundcolor=\color{codebg},
    frame=single,
    rulecolor=\color{codeframe},
    basicstyle=\ttfamily\small,
    showstringspaces=false,
    breaklines=true
}

\title{Laboratorio 9\\Convenciones de Codificación RF.1 y RF.2}
\author{Lorenzo \\ Proyecto ChambeaYa}
\date{\today}

\begin{document}

\maketitle

\section{Datos generales}

\textbf{Requisitos bajo responsabilidad:}

\begin{itemize}
    \item RF.1 - Gestión de Registro y Autenticación.
    \item RF.2 - Gestión del Perfil del Joven.
\end{itemize}

\textbf{Lenguaje y herramientas:}

\begin{itemize}
    \item Python.
    \item Flask.
    \item SQLAlchemy.
    \item VS Code.
    \item SonarLint como apoyo para revisar bugs, code smells y vulnerabilidades.
\end{itemize}

\section{Archivos revisados}

\subsection{RF.1 - Registro y autenticación}

\begin{lstlisting}[style=textstyle]
application/usuario_application_service.py
domain/auth/autenticacion_dominio_servicio.py
domain/auth/i_usuario_repository.py
domain/auth/token_recuperacion.py
domain/auth/usuario.py
frameworks/flask_mvc/routes/auth_routes.py
frameworks/migrations/versions/734d8ba24491_add_user_activation_and_password_reset_.py
frameworks/sqlalchemy_orm/models/usuario_model.py
infrastructure/sqlalchemy_usuario_repository.py
presentation/usuario_controller.py
\end{lstlisting}

\subsection{RF.2 - Gestión del perfil del joven}

\begin{lstlisting}[style=textstyle]
application/perfil_application_service.py
domain/perfil/i_perfil_repository.py
domain/perfil/practicante.py
frameworks/flask_mvc/routes/perfil_routes.py
frameworks/sqlalchemy_orm/models/practicante_model.py
infrastructure/sqlalchemy_perfil_repository.py
presentation/perfil_controller.py
\end{lstlisting}

\section{Convenciones de codificación aplicadas}

\subsection{Práctica 1: nombres en \texttt{snake\_case}}

En Python se recomienda usar \texttt{snake\_case} para nombres de archivos, métodos y variables. Esta convención se aplicó en los módulos de RF.1 y RF.2.

\textbf{Fragmento de código:}

\begin{lstlisting}[style=pythonstyle]
def register_user(self, email, password, tipo):
    self._require_repository()
    normalized_email = self._normalize_email(email)
\end{lstlisting}

\subsection{Práctica 2: clases en \texttt{PascalCase}}

Las clases se nombraron en \texttt{PascalCase}, siguiendo la convención estándar de Python para tipos y clases.

\textbf{Fragmento de código:}

\begin{lstlisting}[style=pythonstyle]
class UsuarioApplicationService:
    def __init__(self, usuario_repository=None, autenticacion_dominio_servicio=None):
        self.usuario_repository = usuario_repository
\end{lstlisting}

\subsection{Práctica 3: separación de responsabilidades por capas}

La implementación respeta la estructura por capas. Las rutas Flask reciben la petición HTTP, los controllers serializan respuestas, los services aplican casos de uso, el dominio contiene reglas de negocio y los repositorios acceden a la base de datos.

\textbf{Fragmento de código:}

\begin{lstlisting}[style=pythonstyle]
def build_usuario_controller():
    repository = SqlAlchemyUsuarioRepository()
    service = UsuarioApplicationService(repository)
    return UsuarioController(service)
\end{lstlisting}

\subsection{Práctica 4: validación temprana de entradas}

Los datos obligatorios se validan al inicio del flujo, evitando continuar con información inválida.

\textbf{Fragmento de código:}

\begin{lstlisting}[style=pythonstyle]
def _normalize_email(self, email):
    if not email:
        raise ValueError("El email es obligatorio")

    return email.strip().lower()
\end{lstlisting}

\subsection{Práctica 5: no exponer objetos internos directamente}

Las respuestas se serializan explícitamente para controlar qué datos salen en JSON. Por ejemplo, no se expone el hash de contraseña.

\textbf{Fragmento de código:}

\begin{lstlisting}[style=pythonstyle]
def _serialize_practicante(self, practicante):
    return {
        "id": practicante.id,
        "usuario_id": practicante.usuario_id,
        "nombres": practicante.nombres,
        "apellidos": practicante.apellidos,
        "habilidades": practicante.habilidades,
    }
\end{lstlisting}

\section{Corrección de bugs, code smells y vulnerabilidades}

\subsection{Corrección 1: eliminar constructores vacíos}

Se eliminaron constructores que solo contenían \texttt{pass}, ya que no aportaban comportamiento y podían ser reportados como code smell.

\textbf{Antes:}

\begin{lstlisting}[style=pythonstyle]
class SqlAlchemyUsuarioRepository(IUsuarioRepository):
    def __init__(self):
        pass
\end{lstlisting}

\textbf{Después:}

\begin{lstlisting}[style=pythonstyle]
class SqlAlchemyUsuarioRepository(IUsuarioRepository):
    def save(self, usuario):
        ...
\end{lstlisting}

\subsection{Corrección 2: reemplazar \texttt{pass} ambiguo}

En métodos fuera del alcance de RF.2 se reemplazó \texttt{pass} por \texttt{NotImplementedError}, dejando explícito que la funcionalidad pertenece a otro requisito.

\textbf{Antes:}

\begin{lstlisting}[style=pythonstyle]
def update_empresa(self):
    pass
\end{lstlisting}

\textbf{Después:}

\begin{lstlisting}[style=pythonstyle]
def update_empresa(self):
    raise NotImplementedError("El perfil de empresa no pertenece a RF.2")
\end{lstlisting}

\subsection{Corrección 3: hashing de contraseñas}

Para evitar vulnerabilidades, las contraseñas no se almacenan en texto plano. Se genera un hash con Werkzeug.

\textbf{Fragmento de código:}

\begin{lstlisting}[style=pythonstyle]
def generar_password_hash(self, password):
    if not password:
        raise ValueError("La contraseña es obligatoria")

    return generate_password_hash(password)
\end{lstlisting}

\subsection{Corrección 4: tokens aleatorios con expiración}

Los tokens de activación y recuperación se generan con valores aleatorios y fecha de expiración.

\textbf{Fragmento de código:}

\begin{lstlisting}[style=pythonstyle]
def generar_token(self, horas_vigencia=24):
    expiracion = datetime.now(timezone.utc) + timedelta(hours=horas_vigencia)
    return TokenRecuperacion(token_urlsafe(32), expiracion)
\end{lstlisting}

\section{Evidencia de avance}

RF.1 y RF.2 fueron integrados previamente en la rama principal del repositorio:

\begin{lstlisting}[style=textstyle]
4ae95d1 feat: implement RF.1 user registration and authentication
51bd587 feat: implement RF.2 joven profile management
\end{lstlisting}

\section{Actualización sugerida en Trello}

\subsection{RF.1}

\textbf{Checklist: Escenarios de Prueba}

\begin{itemize}
    \item Registrar usuario con email, contraseña y tipo válido.
    \item Activar cuenta con token válido.
    \item Iniciar sesión con credenciales válidas.
    \item Rechazar login con contraseña incorrecta.
    \item Recuperar contraseña mediante token.
\end{itemize}

\textbf{Checklist: Tareas de Implementación}

\begin{itemize}
    \item Implementar entidad Usuario.
    \item Implementar servicio de autenticación.
    \item Implementar rutas REST de autenticación.
    \item Implementar repositorio SQLAlchemy.
    \item Validar flujo completo.
\end{itemize}

\subsection{RF.2}

\textbf{Checklist: Escenarios de Prueba}

\begin{itemize}
    \item Crear perfil de practicante autenticado.
    \item Agregar habilidades técnicas.
    \item Agregar formación educativa.
    \item Registrar DNI o carnet universitario.
    \item Consultar score de reputación.
\end{itemize}

\textbf{Checklist: Tareas de Implementación}

\begin{itemize}
    \item Implementar entidad Practicante.
    \item Implementar servicio de perfil.
    \item Implementar rutas REST de perfil.
    \item Implementar repositorio SQLAlchemy.
    \item Validar flujo RF.1 + RF.2.
\end{itemize}

\end{document}
